\section{Finitely presented algebra with proof-relevant equality}

\subsection{Polynomial codes}

Fix a number $n$ of variables.  A monomial is a vector
$\alpha=(\alpha_0,\ldots,\alpha_{n-1})\in\N^n$, encoded by a finite tuple.  A polynomial over
$\Z$ is a finite list of coefficient--monomial pairs.  The code need not be normalized: duplicate
monomials and zero coefficients are allowed.  Evaluation, addition, negation, multiplication,
substitution, coefficient collection on a supplied finite degree box, and equality of collected
coefficient vectors are primitive recursive.

There are two legitimate equality interfaces.

\begin{enumerate}[label=\arabic*.]
\item A canonical normalization algorithm may be used when coefficient equality is decidable and the
      required sorting/collection procedure has been represented.
\item More generally, a coefficientwise identity is supplied as a finite certificate.  This avoids
      requiring a global normal form and extends directly to coefficients in a certified ring.
\end{enumerate}

The second interface is foundational below.  Gröbner bases can accelerate witness production, but
are not required for the definition of quotient equality.

\subsection{Quotient rings by ideal-combination witnesses}

Let
\[
 A=\Z[X_0,\ldots,X_{n-1}]/(g_0,\ldots,g_{r-1}).
\]
An element of $A$ is represented by a polynomial code.  A witness that $p$ and $q$ represent the
same element is a vector of polynomial codes $h=(h_0,\ldots,h_{r-1})$ satisfying
\begin{equation}\label{eq:ideal-witness}
        p-q=\sum_{i<r} h_i g_i
\end{equation}
coefficientwise.

\begin{theorem}[Exact certified presentation of a quotient ring]\label{thm:quotient-setoid}
The checker for \eqref{eq:ideal-witness} and the following witness transformers make polynomial
codes into a certified commutative ring presenting $A$:
\begin{align*}
\mathsf{refl}(p)&=(0,\ldots,0),\\
\mathsf{sym}(h)&=(-h_0,\ldots,-h_{r-1}),\\
\mathsf{trans}(h,k)&=(h_0+k_0,\ldots,h_{r-1}+k_{r-1}).
\end{align*}
Addition, negation, and multiplication admit primitive-recursive congruence transformers.
\end{theorem}

\begin{proof}
Reflexivity, symmetry, and transitivity follow by the corresponding polynomial identities.  Suppose
$p-p'=\sum h_i g_i$ and $q-q'=\sum k_i g_i$.  Then
\[
 (p+q)-(p'+q')=\sum_i(h_i+k_i)g_i,
\]
and
\[
 pq-p'q'=(p-p')q+p'(q-q')
         =\sum_i(h_iq+p'k_i)g_i.
\]
These formulas are the required witness transformers.  Negation sends $h_i$ to $-h_i$.  Ring laws
are literal polynomial identities and therefore have fixed zero or elementary coefficient
certificates.  Every operation is a bounded fold over finite polynomial lists, so its correctness is
proved by quantifier-free induction on a list index.
\end{proof}

The relation is complete by definition of the generated ideal: $p$ and $q$ have equal quotient
classes exactly when some finite list $h_i$ satisfies \eqref{eq:ideal-witness}.  Completeness does not
imply a primitive-recursive search that finds $h$ from $p,q$ in every presentation.  The certificate
may be an input or the output of a separately justified algorithm.

\subsection{Presented algebras over a certified base ring}

The same construction works over any fixed certified commutative ring $R$.  Polynomial coefficients
are $R$-codes; coefficient identities carry $R$-equality witnesses; and the ideal witness includes
both multiplier polynomials and the coefficientwise equality certificates.  A uniform theorem may
quantify over a finite presentation code only when the base-ring operations and their correctness
proofs are supplied through a fixed interface.  It may not call an arbitrary ring program stored in
the input.

Nested presentations can be flattened.  If
$R=\Z[Y]/J$ and $A=R[X]/I$, replace every coefficient by a $Y$-polynomial and present $A$ directly
as a quotient of $\Z[Y,X]$ by $J$ together with lifted generators of $I$.  The flattening algorithm
is primitive recursive and preserves every supplied ideal witness by substitution.

\subsection{Morphisms and their certificates}

Let
\[
 A=\Z[X_0,\ldots,X_{n-1}]/(g_0,\ldots,g_{r-1}),\qquad
 B=\Z[Y_0,\ldots,Y_{m-1}]/(h_0,\ldots,h_{s-1}).
\]
A presented ring map $\phi:A\to B$ consists of polynomials
$p_0(Y),\ldots,p_{n-1}(Y)$ and, for each $g_i$, a witness
\[
  g_i(p_0,\ldots,p_{n-1})=\sum_{j<s} a_{ij}h_j.
\]
The induced map sends $q(X)$ to $q(p(Y))$.  If $q-q'=\sum b_i g_i$, then substitution gives
\[
 q(p)-q'(p)=\sum_i b_i(p)g_i(p)
           =\sum_j\left(\sum_i b_i(p)a_{ij}\right)h_j,
\]
so the map has an explicit equality-witness transformer.

Composition substitutes the image tuple of one map into the other and composes the relation
certificates by the displayed formula.  Identity maps use the coordinate variables.  An isomorphism
certificate is a pair of presented maps plus coordinatewise witnesses that both composites equal the
identities.  Thus affine morphisms and affine isomorphisms are finite arithmetic data.

\subsection{Modules, matrices, kernels, and cokernels}

A finitely presented $A$-module is given by a matrix
$R:A^u\to A^v$ and the cokernel presentation $A^v/\im R$.  A vector equality witness is a coefficient
vector $z\in A^u$ satisfying $x-y=Rz$, with componentwise ring witnesses.  The same witness algebra
as in Theorem~\ref{thm:quotient-setoid} gives a certified module.

Matrix multiplication, direct sums, tensor-presentation matrices, dual maps, Koszul differentials,
and finite chain complexes are all primitive recursive on finite codes.  Kernel membership is
expressed by a checked equation $Rx=0$; image membership by a witness $x=Ry$; and exactness at a
particular element by an explicit preimage witness.  This is exactly the level needed for finite
\v{C}ech arguments.

\subsection{Gröbner bases as optional certificates}

A Gröbner basis is valuable for producing ideal witnesses and normal forms, but the foundational
interface should not assume that a uniform unbounded Buchberger run is a primitive-recursive
oracle.  Three safe modes are available.

\begin{description}[style=nextline]
\item[Verification mode.]
The input contains a proposed basis, representations of its elements in the original ideal,
reduction traces for the original generators, and reductions of all relevant $S$-polynomials to
zero.  A fixed checker verifies Buchberger's criterion and every identity.

\item[Fuelled mode.]
A fixed Buchberger evaluator runs for a supplied bound and returns its finite trace.  Soundness of a
successful output is proved; failure to finish within the bound makes no completeness claim.

\item[Bounded-family mode.]
For a family with a separately proved degree/size bound, the bound becomes the fuel for a total
primitive-recursive constructor.
\end{description}

This separation matters for computational algebraic geometry.  Ideal membership certificates,
Hilbert-function tables, elimination identities, and finite standard-monomial bases can all be
checked in iPRA even when their general-purpose discovery algorithm is not part of the theorem.

\subsection{A small exact example}

In $\Z[x]/(x^2-2)$ the equality $x^4=4$ is certified by
\[
 x^4-4=(x^2+2)(x^2-2).
\]
The executable companion represents each polynomial sparsely and verifies this coefficientwise.
Nothing is inferred from an approximate root of $x^2-2$; the quotient identity is exact finite
algebra.
